\section{The phase at infinity in the light-ring limit}
\label{app:phase}

Use the tortoise origin \(x=0\) at \(r=3\), \(M=1\), and
\(h=(j+1/2)^{-1}\). With \(s=b-1/2\), the semiclassical equation is
\[
 -h^2 u''+(V_0+h^2 V_2)u=E_h u,\quad
 V_2=-f\left(\frac{1}{4r^2}+\frac{6}{r^3}\right),
 \quad E_h=a+2as h+O(h^2).
\]
In the Weber equation the ingoing solution is \(D_\nu(-z)\) with
\(z=e^{-\ii\pi/4}\sqrt{2a}\,y\), \(\nu=-1/2+\ii s\).
Matching the right wave phases \(\pm(ay^2/2+s\log y)\) in the
connection formula of \cite{dlmf} gives
\[
 r_s=-\,\frac{\ii\,\Gamma(1/2-\ii s)}{\sqrt{2\pi}}
      e^{-\pi s/2}(2a)^{\ii s}.
\]
In particular \(|r_s|^2=(1+e^{2\pi s})^{-1}\).
To fix the phase at infinity let
\(p_h=\sqrt{E_h-V_0-h^2 V_2}\) and
\(\Phi_h(x)=\sqrt{E_h}x-\int_x^\infty(p_h(t)-\sqrt{E_h})\,\ud t\).
Write \(p_0=\sqrt{a-V_0}\). The constants that survive right matching are
\[
\begin{split}
 J&=\int_0^\infty(\sqrt a-p_0(t))\,\ud t,\\
 C_E&=\lim_{x\downarrow0}
 \left[\frac{x}{2\sqrt a}
   -\frac12\int_x^\infty\left(\frac1{p_0(t)}-\frac1{\sqrt a}\right)\ud t
   -\frac{\log x}{2a}\right].
\end{split}
\]
Here \(p_0(x)=a x+O(x^2)\) on the right. At \(x=h^{2/5}\) set
\(Y=x/\sqrt h\). The phase divided by \(h\) is
\[
 \tfrac12 aY^2+s\log Y+J/h+\tfrac{s}{2}\log h
    + 2as C_E+O(h^{1/5}).
\]
The static Taylor remainder is \(O(x^3/h)\), and the second-order
energy remainder is \(O(h/x^2)\). The integrals for the retained
terms give \(J\simeq1.09919898\) and \(C_E\simeq-21.23316679\).
Division by the nonzero Weber incident coefficient then gives
\[
 R_j\!\left(\frac{j+b}{3\sqrt3}\right)
   = e^{-2\ii J(j+1/2)+\ii s\log(j+1/2)-4\ii a s C_E}
     \big[r_s+O((j+1/2)^{-1/5})\big].
\]
The factor of two in the phase is from taking the outgoing over
incoming coefficient. On the left the common phase cancels. The
parent and daughter share the same tortoise origin. Their phase
difference yields the coefficient \(1/2\) of \(\log\ell\) in
Section~\ref{sec:followons}.
